\documentclass[conference]{IEEEtran}
\IEEEoverridecommandlockouts
\usepackage{cite}
\usepackage{amsmath,amssymb,amsfonts}
\usepackage{algorithmic}
\usepackage{graphicx}
\usepackage{textcomp}
\usepackage{xcolor}
\usepackage{listings}
\usepackage{hyperref}

\lstset{
  basicstyle=\ttfamily\footnotesize,
  breaklines=true,
  frame=single,
  language=Python,
  showstringspaces=false,
  keywordstyle=\color{blue},
  commentstyle=\color{gray},
  stringstyle=\color{red},
}

\def\BibTeX{{\rm B\kern-.05em{\sc i\kern-.025em b}\kern-.08em
    T\kern-.1667em\lower.7ex\hbox{E}\kern-.125emX}}
\begin{document}

\title{Mauerstrassen Agent}

\author{\IEEEauthorblockN{Aktuelle Themen und Trends I und II \\
BWSE - SS 2026}
\IEEEauthorblockA{\textit{Group 7} \\
Dario-Benjamin Dzakic -- 2310859034\\
Marco Magyar -- 2410859018}
}

\maketitle

\begin{abstract}
Tick the boxes for the level your implementation achieves:

\begin{itemize}
    \item[$\square$] Minimal
    \item[$\square$] Advanced
    \item[$\boxtimes$] Expert
\end{itemize}

\end{abstract}

%----------------------------------------------------------------------
\section{Work Distribution}
%----------------------------------------------------------------------

\textbf{Dario-Benjamin Dzakic} designed and implemented the first three agents:
\emph{Strategy Advisor}, \emph{Market Analyst}, and \emph{Stock Screener},
together with their corresponding tasks (\texttt{strategy\_task},
\texttt{trend\_task}, \texttt{screening\_task}).
He built the sector-analysis RAG knowledge source
(\texttt{sector\_analysis.md}) and authored the Finnhub MCP server
(\texttt{finnhub\_server.py}), which exposes six financial-data tools
(\texttt{get\_company\_profile}, \texttt{get\_basic\_financials},
\texttt{get\_financials\_reported}, \texttt{get\_company\_news},
\texttt{get\_recommendation\_trends}, \texttt{get\_quote}).

\textbf{Marco Magyar} designed and implemented the remaining two agents:
\emph{Stock Checker} and \emph{Portfolio Builder}, together with their
corresponding tasks (\texttt{checker\_task}, \texttt{portfolio\_task}).
He built the investment-strategy and risk-framework RAG knowledge sources
(\texttt{investment\_strategies.md}, \texttt{risk\_frameworks.md}) and
developed the MCP--CrewAI integration layer (\texttt{mcp\_servers.py})
that connects the Finnhub MCP server to the CrewAI tool ecosystem via
\texttt{MCPServerAdapter}.
He also designed the overall crew orchestration, including the
hierarchical process configuration with a manager LLM and the Google
Generative AI embedding provider.

Both team members contributed to architecture design, testing, and
documentation.

%----------------------------------------------------------------------
\section{Problem Statement}
%----------------------------------------------------------------------

Retail investors face a complex, information-rich landscape when
constructing equity portfolios. They must assess their own risk
tolerance, identify current market trends, screen relevant stocks,
validate fundamentals, and finally allocate capital in a diversified
manner. Each of these steps requires different expertise---from
behavioural finance to fundamental analysis---and performing them
manually is time-consuming and error-prone.

\textbf{Mauerstrassen Agent} is a multi-agent system built on
CrewAI \cite{CrewAI2024} that automates this end-to-end portfolio
construction workflow. Given a user's capital, risk tolerance, and
investment horizon, the system produces a concrete, validated portfolio
recommendation with position sizes, buy prices, and stop-loss levels.

\textbf{Assumptions:}
\begin{itemize}
    \item The user has a brokerage account and can execute trades independently.
    \item Market data retrieved via the Finnhub API \cite{Finnhub2024} and
          web searches (Serper) is sufficiently up-to-date for portfolio
          construction.
    \item The system produces \emph{recommendations}, not financial advice
          in a regulatory sense.
\end{itemize}

%----------------------------------------------------------------------
\section{Architecture}
%----------------------------------------------------------------------

The system is implemented as a CrewAI \emph{crew} with five specialised
agents orchestrated by a \textbf{hierarchical process}. A manager LLM
coordinates agent execution, delegates tasks, and synthesises
intermediate results. The architecture integrates three key
technologies:

\begin{enumerate}
    \item \textbf{RAG (Retrieval-Augmented Generation):} Domain knowledge
          files are embedded and made available to agents at inference time
          via CrewAI's \texttt{TextFileKnowledgeSource}.
    \item \textbf{MCP (Model Context Protocol):} A custom Finnhub MCP
          server exposes real-time financial data tools over stdio
          transport, consumed by agents through \texttt{MCPServerAdapter}.
    \item \textbf{External Tools:} The Serper web-search tool provides
          live market information for trend analysis and stock screening.
\end{enumerate}

%----------------------------------------------------------------------
\subsection{Agent Description}
%----------------------------------------------------------------------

The system comprises five agents, each with a distinct role:

\begin{enumerate}

\item \textbf{Strategy Advisor}
(\texttt{strategy\_advisor}) ---
An investment strategy expert that interviews the user via
\texttt{human\_input} to collect capital, risk tolerance, and time
horizon. It validates the user's profile against the RAG knowledge
base (investment strategies, risk frameworks, sector analysis) to
detect contradictions (e.g.\ aggressive risk with a 3-month horizon)
and produces a structured investment strategy with recommended sectors.

\textit{Knowledge sources:} \texttt{strategy\_source} (investment
strategies + risk frameworks), \texttt{sector\_source} (sector
analysis).

\item \textbf{Market Analyst}
(\texttt{market\_analyst}) ---
Searches the web using \texttt{SerperDevTool} for five current market
trends aligned with the user's strategy. Presents trends with
rationales and lets the user select which to pursue via
\texttt{human\_input}.

\textit{Tools:} \texttt{SerperDevTool}.

\item \textbf{Stock Screener}
(\texttt{stock\_screener}) ---
For each selected trend, searches the web for the top five stocks best
positioned to benefit. Filters by market cap and liquidity appropriate
to the risk profile.

\textit{Tools:} \texttt{SerperDevTool}.

\item \textbf{Stock Checker}
(\texttt{stock\_checker}) ---
The validation analyst that uses \textbf{MCP tools} (Finnhub API) to
pull company profiles, financial metrics, reported financials, news,
and analyst recommendations for each candidate stock. Produces a
go/no-go signal per stock.

\textit{Tools:} All six Finnhub MCP tools + \texttt{SerperDevTool}.

\item \textbf{Portfolio Builder}
(\texttt{portfolio\_builder}) ---
Synthesises the validated stock list and the investment strategy into a
final portfolio. Computes allocation percentages, share counts, buy
prices, and stop-loss levels. Outputs a clean markdown table.

\textit{Knowledge sources:} \texttt{strategy\_source}.

\end{enumerate}

%----------------------------------------------------------------------
\subsection{Task Description}
%----------------------------------------------------------------------

Five tasks form the pipeline, each assigned to one agent:

\begin{enumerate}

\item \textbf{Strategy Task} (\texttt{strategy\_task}, agent:
\texttt{strategy\_advisor}) --- Interviews the user and produces a
structured investment strategy (capital, risk level, horizon, strategy
type, 2--3 target sectors). Uses \texttt{human\_input}.

\item \textbf{Trend Task} (\texttt{trend\_task}, agent:
\texttt{market\_analyst}) --- Searches for five market trends matching
the strategy. User selects trends to pursue. Context:
\texttt{strategy\_task}. Uses \texttt{human\_input}.

\item \textbf{Screening Task} (\texttt{screening\_task}, agent:
\texttt{stock\_screener}) --- Identifies up to five stocks per
selected trend. Context: \texttt{strategy\_task},
\texttt{trend\_task}.

\item \textbf{Checker Task} (\texttt{checker\_task}, agent:
\texttt{stock\_checker}) --- Validates each stock candidate using
Finnhub MCP tools and web search. Produces per-stock go/no-go
reports. Context: \texttt{screening\_task}.

\item \textbf{Portfolio Task} (\texttt{portfolio\_task}, agent:
\texttt{portfolio\_builder}) --- Builds the final portfolio as a
markdown table with allocation, shares, prices, and stop-losses.
Output is written to
\texttt{output/portfolio\_recommendation.md}.
Context: \texttt{strategy\_task}, \texttt{checker\_task}.

\end{enumerate}

%----------------------------------------------------------------------
\subsection{Flow}
%----------------------------------------------------------------------

The crew uses a \textbf{hierarchical process}
(\texttt{Process.hierarchical}) where a manager LLM orchestrates the
sequential execution of all five tasks. The interaction flow is:

\begin{enumerate}
    \item The \textbf{manager} delegates the \texttt{strategy\_task}
          to the Strategy Advisor, which interviews the user.
    \item The strategy output feeds into the \texttt{trend\_task},
          delegated to the Market Analyst, which again involves the user.
    \item Both previous outputs flow into the
          \texttt{screening\_task}, delegated to the Stock Screener.
    \item The screening results are passed to the
          \texttt{checker\_task}, where the Stock Checker validates
          each candidate via MCP (Finnhub) and web search.
    \item Finally, the manager delegates the \texttt{portfolio\_task}
          to the Portfolio Builder, which receives the original
          strategy and the validation results to construct the final
          recommendation.
\end{enumerate}

Task context dependencies ensure that each agent receives exactly the
information it needs. The \texttt{human\_input} flag on the first two
tasks creates an interactive loop where the user actively shapes the
investment direction.

%----------------------------------------------------------------------
\section{Minimal}
%----------------------------------------------------------------------

The minimal functionality of the system comprises a working multi-agent
crew with clearly separated roles:

\begin{itemize}
    \item \textbf{Five specialised agents}, each with a distinct role,
          goal, and backstory defined in \texttt{agents.yaml}.
    \item \textbf{Five sequential tasks} defined in \texttt{tasks.yaml},
          with explicit context dependencies ensuring information flows
          correctly through the pipeline.
    \item \textbf{Structured output:} The final portfolio task produces
          a Pydantic-validated \texttt{PortfolioRecommendation} model
          containing stocks, summary, total capital, and risk level.
    \item \textbf{CrewAI orchestration} via the \texttt{@CrewBase}
          decorator pattern, with agents and tasks wired together in
          \texttt{crew.py}.
    \item \textbf{External tool integration:} The \texttt{SerperDevTool}
          gives the Market Analyst and Stock Screener access to real-time
          web search results.
\end{itemize}

%----------------------------------------------------------------------
\section{Advanced}
%----------------------------------------------------------------------

Beyond the minimal system, the following advanced features are
implemented:

\begin{itemize}
    \item \textbf{Human-in-the-loop interaction:} The
          \texttt{strategy\_task} and \texttt{trend\_task} use
          \texttt{human\_input=True}, allowing the user to provide
          investment parameters and select trends interactively.
          This ensures the system adapts to individual user needs
          rather than running fully autonomously.
    \item \textbf{Hierarchical process with manager LLM:} Instead of a
          simple sequential process, the crew uses
          \texttt{Process.hierarchical} with a dedicated manager model
          (\texttt{MANAGER\_MODEL} environment variable). The manager
          coordinates task delegation, handles inter-agent communication,
          and synthesises results---mimicking a real investment team with
          a senior portfolio manager overseeing analysts.
    \item \textbf{Task context chaining:} Tasks declare explicit
          \texttt{context} dependencies (e.g.\ \texttt{portfolio\_task}
          receives context from both \texttt{strategy\_task} and
          \texttt{checker\_task}), enabling rich information flow across
          the pipeline.
    \item \textbf{Markdown output artifact:} The portfolio task writes
          its result to \texttt{output/portfolio\_recommendation.md},
          producing a persistent, human-readable deliverable.
\end{itemize}

%----------------------------------------------------------------------
\section{Expert}
%----------------------------------------------------------------------

The expert-level features distinguish this system through deep
integration of RAG and MCP:

\subsection{RAG --- Retrieval-Augmented Generation}

Three curated knowledge documents are embedded and served to agents at
inference time via CrewAI's \texttt{TextFileKnowledgeSource}:

\begin{itemize}
    \item \texttt{investment\_strategies.md} --- Covers value, growth,
          momentum, dividend, and index/passive strategies with key
          metrics and suitability profiles.
    \item \texttt{risk\_frameworks.md} --- Defines conservative,
          moderate, and aggressive risk levels with allocation limits,
          drawdown tolerances, and common profile contradictions to flag.
    \item \texttt{sector\_analysis.md} --- Analyses seven sectors
          (Technology, Healthcare, Energy, Financials, Consumer
          Discretionary, Consumer Staples, Industrials) with drivers,
          risks, and sector rotation patterns.
\end{itemize}

Two knowledge source objects are constructed:
\texttt{strategy\_source} (investment strategies + risk frameworks) and
\texttt{sector\_source} (sector analysis). These are injected into the
\texttt{strategy\_advisor} (both sources) and the
\texttt{portfolio\_builder} (\texttt{strategy\_source}), giving these
agents grounded domain knowledge that goes beyond what the base LLM
knows.

Embeddings are generated using Google Generative AI
(\texttt{google-generativeai}) configured via environment variables,
enabling semantic retrieval over the knowledge base.

\subsection{MCP --- Model Context Protocol}

A custom \textbf{Finnhub MCP server}
(\texttt{finnhub\_server.py}) is implemented using the
\texttt{FastMCP} framework. It exposes six financial-data tools over
\textbf{stdio transport}:

\begin{enumerate}
    \item \texttt{get\_company\_profile} --- Company name, sector,
          industry, market cap, exchange.
    \item \texttt{get\_basic\_financials} --- P/E, P/B, ROE, margins,
          52-week range, dividend yield.
    \item \texttt{get\_financials\_reported} --- As-reported income
          statement, balance sheet, cash flow (annual/quarterly).
    \item \texttt{get\_company\_news} --- Recent company news articles.
    \item \texttt{get\_recommendation\_trends} --- Analyst buy/hold/sell
          recommendation counts.
    \item \texttt{get\_quote} --- Real-time stock quote (price, change,
          high/low, open, previous close).
\end{enumerate}

The MCP server is integrated into CrewAI through
\texttt{MCPServerAdapter} in \texttt{mcp\_servers.py}, which spawns the
server as a subprocess via \texttt{StdioServerParameters} and converts
MCP tools into CrewAI-compatible tool objects. The
\texttt{stock\_checker} agent receives all six MCP tools alongside the
web search tool, enabling it to pull real financial data for stock
validation.

This architecture cleanly separates the data-access layer (MCP server)
from the agent logic (CrewAI), following the Model Context Protocol
standard for tool interoperability.

%----------------------------------------------------------------------
\section{Reflection}
%----------------------------------------------------------------------

During development and testing we encountered several issues that
highlight the practical challenges of building LLM-based multi-agent
systems.

\textbf{MCP context size of Finnhub responses.}
The Finnhub API returns large JSON payloads---especially
\texttt{get\_financials\_reported}, which includes full income
statements, balance sheets, and cash-flow data across multiple periods.
When the \texttt{stock\_checker} agent called this tool for several
stocks in sequence, the accumulated context quickly grew very large.
This caused the LLM to hit token limits or to lose track of earlier
validation results, degrading the quality of the go/no-go assessments.
A possible mitigation would be to add a summarisation step inside the
MCP tool or to limit the number of reporting periods returned.

\textbf{Human input requested but not declared on the crew.}
The \texttt{strategy\_task} and \texttt{trend\_task} set
\texttt{human\_input=True} at the task level, which prompts the user
for input during execution. However, the crew itself was not explicitly
configured to expect human interaction, which led to unexpected
blocking behaviour in certain execution environments. Ensuring that
the crew-level configuration and the task-level flags are aligned is
an important lesson for future projects.

\textbf{Stock count not consistently five.}
The \texttt{screening\_task} instructs the Stock Screener to return
the ``top 5 stocks'' per selected trend. In practice, the agent
sometimes returned fewer or more than five---depending on how the LLM
interpreted the prompt and the availability of suitable candidates in
the search results. Strict enforcement of output cardinality is
difficult with free-form LLM generation; a post-processing validation
step or a stricter output schema could help enforce the constraint.

\textbf{Structured output not always followed by Portfolio Builder.}
A Pydantic model (\texttt{PortfolioRecommendation}) was defined to
structure the final output, but the Portfolio Builder did not
consistently produce output that matched the schema. Fields were
sometimes missing or formatted differently (e.g.\ allocation as a
string instead of a number). This is a known limitation of relying on
LLMs to produce structured output without explicit function-calling
or JSON-mode enforcement---the agent sometimes preferred to write
free-form markdown instead of adhering to the schema.

\end{document}
